% Preámbulo compartido por todos los documentos del proyecto (compilar con pdflatex).
\usepackage[utf8]{inputenc}
\usepackage[T1]{fontenc}
\usepackage{lmodern}
\usepackage{textcomp}
\usepackage[spanish,es-nodecimaldot,es-noshorthands,es-tabla]{babel}
\DeclareUnicodeCharacter{03A9}{\ensuremath{\Omega}}
\DeclareUnicodeCharacter{00B5}{\ensuremath{\mu}}
\DeclareUnicodeCharacter{03BC}{\ensuremath{\mu}}
\DeclareUnicodeCharacter{221A}{\ensuremath{\surd}}
\DeclareUnicodeCharacter{2192}{\ensuremath{\rightarrow}}
\DeclareUnicodeCharacter{2248}{\ensuremath{\approx}}
\DeclareUnicodeCharacter{2212}{\ensuremath{-}}
\DeclareUnicodeCharacter{2264}{\ensuremath{\le}}
\DeclareUnicodeCharacter{2265}{\ensuremath{\ge}}
\DeclareUnicodeCharacter{00D7}{\ensuremath{\times}}
\usepackage[a4paper,top=2.5cm,bottom=2.5cm,left=2.4cm,right=2.4cm,headheight=15pt]{geometry}
\setlength{\emergencystretch}{3em}
\usepackage{amsmath,amssymb,mathtools,bm}
\usepackage{siunitx}
\sisetup{output-decimal-marker={.},group-digits=false,per-mode=symbol,separate-uncertainty}
\usepackage{graphicx,float,booktabs,tabularx,longtable,array,multirow}
\usepackage[font=small,labelfont=bf]{caption}
\usepackage{enumitem}
\setlist{itemsep=2pt,topsep=4pt}
\usepackage[dvipsnames,table]{xcolor}
\usepackage[most]{tcolorbox}
\usepackage{tikz}
\usetikzlibrary{arrows.meta,positioning,shapes.geometric,calc,fit,backgrounds,decorations.markings}
\usepackage[european,straightvoltages]{circuitikz}
\usepackage{listings}
\usepackage{fancyhdr}
\usepackage{titlesec}
\usepackage{xurl}
\usepackage[colorlinks=true,linkcolor=Azul,urlcolor=Azul,citecolor=Azul]{hyperref}

\graphicspath{{../../simulacion/graficas/}{figuras/}}

% ---------------------------------------------------------------- colores
\definecolor{Azul}{HTML}{0D47A1}
\definecolor{Cian}{HTML}{00838F}
\definecolor{Naranja}{HTML}{E65100}
\definecolor{Rojo}{HTML}{C62828}
\definecolor{Verde}{HTML}{2E7D32}
\definecolor{Gris}{HTML}{455A64}
\definecolor{FondoCodigo}{HTML}{F5F7FA}

% ---------------------------------------------------------------- títulos
\titleformat{\section}{\Large\bfseries\color{Azul}}{\thesection}{0.8em}{}
\titleformat{\subsection}{\large\bfseries\color{Cian}}{\thesubsection}{0.8em}{}
\titleformat{\subsubsection}{\normalsize\bfseries\color{Gris}}{\thesubsubsection}{0.8em}{}

% ---------------------------------------------------------------- cabeceras
\pagestyle{fancy}
\fancyhf{}
\fancyhead[L]{\small\color{Gris}\nombredocumento}
\fancyhead[R]{\small\color{Gris}Generador de flujo axial con imanes NdFeB}
\fancyfoot[C]{\small\thepage}
\renewcommand{\headrulewidth}{0.4pt}

% ---------------------------------------------------------------- cajas
\newtcolorbox{nota}[1][Nota]{enhanced,breakable,colback=Azul!4,colframe=Azul!70,
  fonttitle=\bfseries,title={#1},left=6pt,right=6pt}
\newtcolorbox{advertencia}[1][Advertencia]{enhanced,breakable,colback=Rojo!4,colframe=Rojo!80,
  fonttitle=\bfseries,title={#1},left=6pt,right=6pt}
\newtcolorbox{resultado}[1][Resultado]{enhanced,breakable,colback=Verde!5,colframe=Verde!75,
  fonttitle=\bfseries,title={#1},left=6pt,right=6pt}
\newtcolorbox{paso}[1]{enhanced,breakable,colback=white,colframe=Cian!80,
  fonttitle=\bfseries,title={#1},left=6pt,right=6pt}

% ---------------------------------------------------------------- código
\lstdefinestyle{base}{
  basicstyle=\ttfamily\footnotesize,
  backgroundcolor=\color{FondoCodigo},
  frame=single,rulecolor=\color{Gris!40},
  breaklines=true,columns=fullflexible,keepspaces=true,
  keywordstyle=\color{Azul}\bfseries,
  commentstyle=\color{Verde}\itshape,
  stringstyle=\color{Naranja},
  showstringspaces=false,
  extendedchars=true,
  literate={á}{{\'a}}1 {é}{{\'e}}1 {í}{{\'i}}1 {ó}{{\'o}}1 {ú}{{\'u}}1 {ñ}{{\~n}}1
           {Á}{{\'A}}1 {É}{{\'E}}1 {Í}{{\'I}}1 {Ó}{{\'O}}1 {Ú}{{\'U}}1 {Ñ}{{\~N}}1
           {·}{{$\cdot$}}1 {→}{{$\rightarrow$}}1 {Ø}{{\O}}1 {Ω}{{$\Omega$}}1 {µ}{{$\mu$}}1
           {−}{{-}}1 {²}{{$^2$}}1 {√}{{$\surd$}}1 {≈}{{$\approx$}}1
}
\lstdefinestyle{python}{style=base,language=Python}
\lstdefinelanguage{JavaScript}{
  keywords={const,let,var,function,return,class,constructor,new,if,else,for,while,this,true,false,null,typeof,module,export,import},
  sensitive=true,comment=[l]{//},morecomment=[s]{/*}{*/},morestring=[b]',morestring=[b]",morestring=[b]`
}
\lstdefinestyle{js}{style=base,language=JavaScript}
\lstdefinestyle{consola}{style=base,keywordstyle=\color{black}}

% ---------------------------------------------------------------- utilidades
\newcommand{\vect}[1]{\mathbf{#1}}
\newcommand{\unitario}[1]{\mathbf{a}_{#1}}
\newcommand{\dif}{\mathrm{d}}
\DeclareUrlCommand\archivo{\urlstyle{tt}}
\makeatletter
\newcommand{\filasdesde}[1]{\@@input #1 }
\makeatother

\newcommand{\portada}[3]{%
  \begin{titlepage}
    \begin{tikzpicture}[remember picture,overlay]
      \fill[Azul] (current page.north west) rectangle ([yshift=-7cm]current page.north east);
      \node[anchor=north west,text=white,font=\sffamily\bfseries\Huge,text width=16cm,align=left]
        at ([xshift=2.4cm,yshift=-2.2cm]current page.north west) {#1};
      \node[anchor=north west,text=white!85!Azul,font=\sffamily\Large,text width=16cm,align=left]
        at ([xshift=2.4cm,yshift=-4.6cm]current page.north west) {#2};
    \end{tikzpicture}
    \vspace*{7cm}
    \begin{center}
      \begin{tikzpicture}[scale=0.9]
        \foreach \i in {0,...,7} {
          \pgfmathsetmacro{\ang}{45*\i}
          \ifodd\i \def\col{Azul!70}\def\pol{S}\else \def\col{Rojo!80}\def\pol{N}\fi
          \fill[\col] (\ang:2.2) circle (0.55);
          \node[white,font=\bfseries] at (\ang:2.2) {\pol};
        }
        \draw[Gris,line width=1pt] (0,0) circle (3.0);
        \fill[Gris!40] (0,0) circle (0.35);
      \end{tikzpicture}
    \end{center}
    \vfill
    {\large #3\par}
    \vspace{0.6cm}
    {\small\color{Gris}Basado en M.\,N.\,O. Sadiku, \emph{Elementos de Electromagnetismo}, 3.ª ed., Oxford University Press.\par
    Valores numéricos generados por el paquete \texttt{generador\_axial} (\texttt{uv run generador-axial reporte}).\par}
  \end{titlepage}
}

% Archivo generado automáticamente por `generador-axial reporte`. No editar a mano.
\newcommand{\ValGradoIman}{N38}  % Grado comercial del imán
\newcommand{\ValBr}{1.25}  % Remanencia [T]
\newcommand{\ValDiametroIman}{20}  % [mm]
\newcommand{\ValRadioIman}{10.0}  % [mm]
\newcommand{\ValEspesorIman}{5.0}  % [mm]
\newcommand{\ValAreaImanMM}{314.16}  % [mm^2]
\newcommand{\ValPolos}{8}
\newcommand{\ValParesPolos}{4}
\newcommand{\ValPasoPolarGrados}{45.0}  % [grados]
\newcommand{\ValRadioMontaje}{40.0}  % [mm]
\newcommand{\ValBobinas}{8}
\newcommand{\ValEspiras}{180}
\newcommand{\ValEspirasTotales}{1440}
\newcommand{\ValDiametroInteriorBobina}{12.0}  % [mm]
\newcommand{\ValDiametroExteriorBobina}{28.0}  % [mm]
\newcommand{\ValAlturaBobina}{6.0}  % [mm]
\newcommand{\ValEntrehierro}{2.5}  % [mm]
\newcommand{\ValEntrehierroMasAltura}{8.5}  % [mm]
\newcommand{\ValAcoplamiento}{0.85}
\newcommand{\ValCalibreAWG}{26}
\newcommand{\ValDiametroCobre}{0.404}  % [mm]
\newcommand{\ValDiametroEsmaltado}{0.455}  % [mm]
\newcommand{\ValAreaCobreMM}{0.1282}  % [mm^2]
\newcommand{\ValResistividadCobre}{1.72}  % [1e-8 ohm m]
\newcommand{\ValLongitudMediaEspira}{62.83}  % [mm]
\newcommand{\ValLongitudAlambreBobina}{11.31}  % [m]
\newcommand{\ValLongitudAlambreTotal}{90.48}  % [m]
\newcommand{\ValResistenciaBobina}{1.517}  % [ohm]
\newcommand{\ValResistenciaInterna}{12.14}  % [ohm]
\newcommand{\ValVentanaBobinaMM}{48.0}  % [mm^2]
\newcommand{\ValOcupacion}{0.78}
\newcommand{\ValRaizCaraIman}{11.1803}  % sqrt(L^2+R^2) [mm]
\newcommand{\ValTerminoCaraIman}{0.4472}  % L/sqrt(L^2+R^2)
\newcommand{\ValCampoCaraIman}{0.2795}  % [T]
\newcommand{\ValEntrehierroMasEspesor}{7.5}  % g+L [mm]
\newcommand{\ValRaizGL}{12.5000}  % sqrt((g+L)^2+R^2) [mm]
\newcommand{\ValRaizG}{10.3078}  % sqrt(g^2+R^2) [mm]
\newcommand{\ValTerminoUnoG}{0.6000}  % (g+L)/sqrt((g+L)^2+R^2)
\newcommand{\ValTerminoDosG}{0.2425}  % g/sqrt(g^2+R^2)
\newcommand{\ValCampoCaraBobina}{0.2234}  % [T]
\newcommand{\ValEntrehierroAlturaEspesor}{13.5}  % g+h+L [mm]
\newcommand{\ValRaizGHL}{16.8003}  % sqrt((g+h+L)^2+R^2) [mm]
\newcommand{\ValRaizGH}{13.1244}  % sqrt((g+h)^2+R^2) [mm]
\newcommand{\ValPrimitivaEnG}{2.1922}  % F(g) [mm]
\newcommand{\ValPrimitivaEnGH}{3.6759}  % F(g+h) [mm]
\newcommand{\ValDiferenciaPrimitivas}{1.4837}  % F(g+h)-F(g) [mm]
\newcommand{\ValCampoMedioBobina}{0.1545}  % [T]
\newcommand{\ValReduccionPorPromedio}{30.8}  % [%]
\newcommand{\ValFlujoMaximoMicro}{41.270}  % [uWb]
\newcommand{\ValAreaImanM}{3.1416}  % [1e-4 m^2]
\newcommand{\ValCaidaDiodo}{0.30}  % [V]
\newcommand{\ValCaidaPuente}{0.60}  % [V]
\newcommand{\ValCapacitanciaMicro}{220}  % [uF]
\newcommand{\ValTensionNominalCondensador}{50}  % [V]
\newcommand{\ValRpmMaxima}{1200}
\newcommand{\ValTensionDcMaxima}{29.27}  % [V]
\newcommand{\ValMargenCondensador}{41}  % [%]
\newcommand{\ValResistenciaLed}{2200}  % [ohm]
\newcommand{\ValPotenciaNominalResistenciaLed}{0.50}  % [W]
\newcommand{\ValCaidaLed}{2.0}  % [V]
\newcommand{\ValCorrienteLedRefMili}{2.21}  % [mA]
\newcommand{\ValCorrienteLedMaxMili}{12.4}  % [mA]
\newcommand{\ValCorrienteLedLimiteMili}{20}  % [mA]
\newcommand{\ValPotenciaResistenciaLedMax}{0.338}  % [W]
\newcommand{\ValRizadoRef}{251.4}  % [mV]
\newcommand{\ValRpmReferencia}{300}
\newcommand{\ValOmegaMecanicaRef}{31.416}  % [rad/s]
\newcommand{\ValFrecuenciaRef}{20.00}  % [Hz]
\newcommand{\ValPeriodoRefMs}{50.00}  % [ms]
\newcommand{\ValOmegaElectricaRef}{125.66}  % [rad/s]
\newcommand{\ValTensionPicoRef}{7.47}  % [V]
\newcommand{\ValTensionRmsRef}{5.28}  % [V]
\newcommand{\ValTensionDcRef}{6.87}  % [V]
\newcommand{\ValTensionPicoPorBobinaRef}{0.933}  % [V]
\newcommand{\ValRpmCarga}{400}
\newcommand{\ValOmegaMecanicaCarga}{41.89}  % [rad/s]
\newcommand{\ValTensionPicoCarga}{9.96}  % [V]
\newcommand{\ValTensionRmsCarga}{7.04}  % [V]
\newcommand{\ValCorrienteCarga}{0.290}  % [A]
\newcommand{\ValTensionEnCarga}{3.52}  % [V]
\newcommand{\ValPotenciaCarga}{1.021}  % [W]
\newcommand{\ValPotenciaPerdidas}{1.021}  % [W]
\newcommand{\ValPotenciaMecanica}{2.042}  % [W]
\newcommand{\ValParResistenteMili}{48.7}  % [mN m]
\newcommand{\ValEficiencia}{50}  % [%]
\newcommand{\ValBateriaTension}{12}  % [V]
\newcommand{\ValBateriaCapacidad}{7}  % [Ah]
\newcommand{\ValBateriaEnergia}{84}  % [Wh]
\newcommand{\ValBateriaResistencia}{0.05}  % [ohm]
\newcommand{\ValBateriaVacia}{11.8}  % E_oc con SoC = 0 [V]
\newcommand{\ValBateriaLlena}{12.8}  % E_oc con SoC = 1 [V]
\newcommand{\ValBateriaPendiente}{1.0}  % [V]
\newcommand{\ValSocInicial}{50}  % [%]
\newcommand{\ValSocCorte}{20}  % [%]
\newcommand{\ValSocReconexion}{25}  % [%]
\newcommand{\ValTensionOcRef}{12.30}  % [V]
\newcommand{\ValEnergiaUtilRef}{25.2}  % [Wh]
\newcommand{\ValControladorEficiencia}{90}  % [%]
\newcommand{\ValControladorArranque}{3.5}  % [V]
\newcommand{\ValControladorCorrienteMax}{3.0}  % [A]
\newcommand{\ValControladorAbsorcion}{14.4}  % [V]
\newcommand{\ValRpmMicrored}{600}
\newcommand{\ValFemRmsMicrored}{10.56}  % [V]
\newcommand{\ValTensionDcMicrored}{14.34}  % [V]
\newcommand{\ValPotenciaExtraibleMicrored}{2.297}  % [W]
\newcommand{\ValPotenciaHaciaBateria}{2.067}  % [W]
\newcommand{\ValCorrienteCargaMicroredMili}{168}  % [mA]
\newcommand{\ValPotenciaMecanicaMicrored}{4.59}  % [W]
\newcommand{\ValParMicroredMili}{73.1}  % [mN m]
\newcommand{\ValHorasRecarga}{20.8}  % [h]
\newcommand{\ValInversorTension}{110}  % [V]
\newcommand{\ValInversorPico}{155.6}  % [V]
\newcommand{\ValInversorTensionAlta}{220}  % [V]
\newcommand{\ValInversorPicoAlto}{311.1}  % [V]
\newcommand{\ValInversorFrecuencia}{60}  % [Hz]
\newcommand{\ValInversorNominal}{150}  % [W]
\newcommand{\ValInversorEficiencia}{90}  % [%]
\newcommand{\ValRelacionEspiras}{13.0}  % V_pico,AC / V_bat
\newcommand{\ValPotenciaCargaAc}{15}  % [W]
\newcommand{\ValCorrienteAcMili}{136}  % [mA]
\newcommand{\ValPotenciaEntradaInversor}{16.67}  % [W]
\newcommand{\ValCorrienteInversorBateria}{1.36}  % [A]
\newcommand{\ValPotenciaNetaBateria}{-14.60}  % [W]
\newcommand{\ValAutonomiaHoras}{1.73}  % [h]
\newcommand{\ValAutonomiaParadoHoras}{1.51}  % [h]
\newcommand{\ValRpmEquilibrio}{1704}
\newcommand{\ValPotenciaCargaReducida}{5}  % [W]
\newcommand{\ValRpmEquilibrioReducida}{984}
\newcommand{\ValInerciaRotor}{2.415}  % [1e-4 kg m^2]
\newcommand{\ValInductanciaBobinaMili}{0.528}  % [mH]
\newcommand{\ValInductanciaTotalMili}{4.22}  % [mH]
\newcommand{\ValFlujoMaximoReal}{0.827}  % phi_max normalizado
\newcommand{\ValPicoDerivadaFlujo}{0.746}
\newcommand{\ValRmsDerivadaFlujo}{0.559}
\newcommand{\ValFactorCresta}{1.335}
\newcommand{\ValFactorMedia}{0.706}
\newcommand{\ValRelacionFemRealIdeal}{0.930}
\newcommand{\ValPuntosRadiales}{32}
\newcommand{\ValMuestrasPerfil}{256}
\newcommand{\ValPasoIntegracionMili}{5}  % [ms]
\newcommand{\ValParCoulombMili}{3.0}  % [mN m]
\newcommand{\ValCoefViscosoMicro}{15.0}  % [uN m s]
\newcommand{\ValCoefVentilacionNano}{15.0}  % [nN m s^2]
\newcommand{\ValParMaximoMotriz}{0.30}  % [N m]
\newcommand{\ValVelocidadVacioMotriz}{1500}  % [RPM]
\newcommand{\ValParImpulso}{0.08}  % [N m]
\newcommand{\ValDuracionImpulso}{0.5}  % [s]
\newcommand{\ValCapacidadTermica}{100}  % [J/K]
\newcommand{\ValConveccionNatural}{8}  % [W/m^2K]
\newcommand{\ValConveccionForzada}{0.15}  % [W/m^2K por rad/s]
\newcommand{\ValCoefCobre}{0.393}  % [%/K]
\newcommand{\ValCoefRemanencia}{-0.12}  % [%/K]
\newcommand{\ValTemperaturaMaximaPla}{55}  % [C]
\newcommand{\ValBateriaRCero}{25}  % [mohm]
\newcommand{\ValBateriaRPolarizacion}{30}  % [mohm]
\newcommand{\ValBateriaTau}{30}  % [s]
\newcommand{\ValBateriaPeukert}{1.12}
\newcommand{\ValBateriaEtaCoulomb}{95}  % [%]
\newcommand{\ValBateriaEtaLlena}{60}  % [%]
\newcommand{\ValBateriaSocGas}{85}  % [%]
\newcommand{\ValBateriaAutodescarga}{3}  % [%/mes]
\newcommand{\ValBateriaSocFlotacion}{99.5}  % [%]
\newcommand{\ValBateriaFactorGas}{1500}
\newcommand{\ValBateriaPendienteGas}{50}
\newcommand{\ValControladorFlotacion}{13.8}  % [V]
\newcommand{\ValConsumoControladorMili}{8}  % [mA]
\newcommand{\ValInversorVacio}{3.6}  % [W]
\newcommand{\ValInversorCoefCuadratico}{5.81}  % [1e-4 1/W]
\newcommand{\ValInversorCorte}{10.8}  % [V]
\newcommand{\ValInversorReconexion}{12.3}  % [V]
\newcommand{\ValParReferencia}{0.07}  % [N m]
\newcommand{\ValRpmRegimenRef}{483}
\newcommand{\ValPotenciaEjeRef}{2.40}  % [W]
\newcommand{\ValPotenciaElectromagneticaRef}{2.21}  % [W]
\newcommand{\ValPerdidasCobreRef}{0.986}  % [W]
\newcommand{\ValPerdidasRozamientoRef}{0.192}  % [W]
\newcommand{\ValPerdidasFoucaultMili}{1.65}  % [mW]
\newcommand{\ValPotenciaBateriaRef}{0.96}  % [W]
\newcommand{\ValEficienciaRef}{40}  % [%]
\newcommand{\ValTensionDcRefReal}{3.50}  % [V]
\newcommand{\ValCorrienteGeneradorRef}{0.281}  % [A]
\newcommand{\ValTemperaturaRegimenRef}{27.8}  % [C]
\newcommand{\ValRpmConCargaRef}{483}
\newcommand{\ValPotenciaEntradaInversorRef}{18.7}  % [W]
\newcommand{\ValCorrienteNetaConCarga}{-1.48}  % [A]
\newcommand{\ValRpmParMaximo}{1002}
\newcommand{\ValPotenciaBateriaParMaximo}{4.35}  % [W]
\newcommand{\ValRizadoBajo}{18.8}  % [V] con 220 uF
\newcommand{\ValRizadoAlto}{0.88}  % [V] con 4700 uF
\newcommand{\ValRpmPicoImpulso}{505}
\newcommand{\ValTiempoParadaImpulso}{2.7}  % [s]
\newcommand{\ValHorasHastaAbsorcion}{4.0}  % [h]
\newcommand{\ValHorasHastaFlotacion}{8.2}  % [h]
\newcommand{\ValCorrienteFlotacionMili}{21}  % [mA]

